% TOP_PROBLEM: {"id": 10000043, "title": "Infinite-cluster intersections with vertical fibers", "category_id": 19, "category": {"id": 19, "name": "probability", "display_name": "Probability", "description": "Problems involving probability theory, stochastic processes, and random structures.", "slug": "probability", "order_index": 19, "created_at": "2026-07-31T15:26:25.671Z"}, "status": "open", "problem_number": "AMR-099-0043", "set_id": 13}
% FIRST_POSTED: 2026-10-01
% ENTRY_KIND: statement_only
% UnsolvedMath ID 10000043; source catalogue code AMR-099-0043
% !TeX program = lualatex
% Definition and sources only; this file is not an LLM solution attempt.
\documentclass[11pt,a4paper]{article}
\usepackage[margin=25mm]{geometry}
\usepackage{fontspec}
\setmainfont{lmroman10-regular.otf}[BoldFont=lmroman10-bold.otf,
 ItalicFont=lmroman10-italic.otf,BoldItalicFont=lmroman10-bolditalic.otf]
\usepackage{amsmath,amssymb,mathtools}
\usepackage{unicode-math}
\setmathfont{latinmodern-math.otf}
\AtBeginDocument{\let\setminus\smallsetminus}
\usepackage{xurl,hyperref}
\hypersetup{colorlinks=true,urlcolor=blue}
\setcounter{secnumdepth}{0}
\setlength{\parindent}{0pt}
\setlength{\parskip}{5pt}
\setlength{\emergencystretch}{3em}
\begin{document}
\section*{Infinite-cluster intersections with vertical fibers}\label{10000043}
{\small UnsolvedMath ID 10000043 · AMR-099-0043 · Probability · AMR Open Problem Lists}
\subsection{Definitions and mathematical statement}
Let $G$ be an infinite connected locally finite graph. For independent bond percolation on $G$, let
\[
p_c(G)=\inf\{q\in[0,1]:\mathbb P_q(o\text{ lies in an infinite open component})>0\},
\]
where $o$ is any vertex. Suppose $p_c(G)=1$, so no constant parameter below one produces an infinite component.

Form the Cartesian product $G\times\mathbb Z$. It has vertices $(v,n)$ and edges $(v,n)\sim(w,n)$ when $v\sim_Gw$, together with $(v,n)\sim(v,n+1)$. For a fixed $p\in[0,1]$, retain all product edges independently with probability $p$. The vertical fiber above $v$ is $F_v=\{v\}\times\mathbb Z$.

Is it almost surely true that every infinite open cluster $C$ and every vertex $v\in V(G)$ satisfy
\[
C\cap F_v\ne\varnothing\quad\Longrightarrow\quad |C\cap F_v|=\infty?
\]
The quantification is simultaneous over all clusters and all fibers in one realization, for each fixed parameter. If no infinite cluster exists, the assertion is vacuous; at $p=1$, the entire connected product is one cluster.

A cluster meeting a fiber infinitely often need not contain any infinite vertical run of open edges in that fiber. Its repeated returns may be linked through arbitrarily distant horizontal excursions. The hypothesis concerns the critical probability of $G$ itself, whereas the random cluster lives in its product with the line, which can have a smaller threshold. No transitivity of $G$ is imposed.
\subsection{Short English statement}
When the base graph cannot percolate below parameter one, must every infinite cluster in its product with the line return infinitely often to each fiber it visits?
\subsection{Sources}
\begin{itemize}
\item[S1] ULAM.AI and the UnsolvedMath Contributors, supplied problems.json, record 10000043 (AMR-099-0043), AMR Open Problem Lists. Catalogue identity and source transcription; CC BY 4.0.
\url{https://huggingface.co/datasets/ulamai/UnsolvedMath}.
\item[S2] Benjamini - Coarse Geometry and Randomness (Saint-Flour notes), Open problem 9.49, PDF page 76. Original source indexed by AMR-099-0043.
\url{https://www.wisdom.weizmann.ac.il/~itai/stflouraug24.pdf#page=76}.
\item[S3] I. Benjamini, Coarse Geometry and Randomness, primary Saint-Flour notes; the numbered question and its surrounding definitions.
\url{https://arquivo.pt/noFrame/replay/20201231041548id_/http://www.wisdom.weizmann.ac.il/~itai/stflouraug24.pdf}.
\end{itemize}
\end{document}
